\subsectionaddtolist{روش‌های ترکیب \lr{MAC} و رمزنگاری}

\begin{enumerate}
    \item {شرح دو روش بر اساس عملیات روی \lr{Plaintext} و \lr{Ciphertext}:}
    \begin{itemize}
        \item {روش اول \lr{MAC-then-Encrypt - MtE}:} در این روش ابتدا تگ \lr{MAC} روی متن اصلی ($P$) محاسبه می‌شود. سپس کل مجموع شامل متن اصلی و تگ \lr{MAC} با هم رمزنگاری می‌شوند تا متن رمزشده نهایی ($C$) تولید شود:
        $$C = E_{k_1}(P \parallel \text{\lr{MAC}}_{k_2}(P))$$
        \item {روش دوم \lr{Encrypt-then-MAC - EtM}:} در این روش ابتدا متن اصلی ($P$) رمزنگاری می‌شود تا متن رمزشده ($C$) به دست آید. سپس تگ \lr{MAC} روی این متن رمزشده محاسبه و به آن الحاق می‌گردد:
        $$C = E_{k_1}(P), \quad t = \text{\lr{MAC}}_{k_2}(C) \implies \text{Result} = C \parallel t$$
    \end{itemize}

    \item {روش توصیه‌شده در علم رمزنگاری نوین:}
    روش \lr{Encrypt-then-MAC (EtM)} عموما توصیه می‌شود. 
    این روش امنیت در برابر حملات متن رمزشده انتخابی (CCA) را به طور کامل تضمین می‌کند. گیرنده می‌تواند پیش از انجام عملیات رمزگشایی (که از نظر محاسباتی سنگین است و می‌تواند آسیب‌پذیری‌های امنیتی ایجاد کند)، اصالت و یکپارچگی بسته را بررسی کرده و در صورت دستکاری، آن را فورا رد کند.

    \item {سناریوی حمله (حمله \lr{Padding Oracle Attack}):}
    در روش \lr{MAC-then-Encrypt}، اگر مهاجم متن رمزشده را دستکاری کند، سیستم ابتدا مجبور است آن را رمزگشایی کند تا به تگ \lr{MAC} دسترسی پیدا کند و صحت آن را بسنجد. در حین رمزگشایی، اگر \lr{Padding} پیام نادرست باشد، سیستم ممکن است خطای متفاوتی نسبت به زمان نادرست بودن \lr{MAC} صادر کند (یا رفتار زمانی متفاوتی نشان دهد). مهاجم با بهره‌برداری از این تفاوت رفتار \lr{(Oracle)}، می‌تواند بایت به بایت متن اصلی را بدون داشتن کلید کشف کند. در حالی که در روش \lr{EtM}، پیام دستکاری‌شده قبل از هرگونه رمزگشایی به دلیل نامعتبر بودن \lr{MAC} رد می‌شود و این حمله عملا غیرممکن می‌گردد.
\end{enumerate}

\subsectionaddtolist{محدوده اعمال \lr{MAC} روی بسته‌های شبکه}

\begin{enumerate}
    \item {انتخاب با امنیت بیشتر:}
    اعمال \lr{MAC} روی {ترکیب هدر و بارمفید $(h, p)$} امنیت به مراتب بیشتری را فراهم می‌کند؛ زیرا یکپارچگی کل بسته (هم داده و هم اطلاعات کنترلی) حفظ می‌شود.

    \item {توصیف حمله در صورت عدم لحاظ هدر ($h$):}
    {حمله تغییر هدر (\lr{Header Modification / Spoofing}):} اگر \lr{MAC} فقط روی $p$ اعمال شود، مهاجم در طول مسیر می‌تواند فیلدهای مهم هدر (مانند آدرس آی‌پی مبدأ/مقصد، شماره پورت‌ها یا فیلدهای اولویت) را دستکاری کند. از آنجا که بارمفید تغییر نکرده، تگ \lr{MAC} همچنان معتبر باقی می‌ماند و گیرنده متوجه دستکاری هدر نخواهد شد. این امر می‌تواند منجر به هدایت اشتباه ترافیک، دور زدن فایروال‌ها یا حملات منع سرویس (\lr{DoS}) شود.

    \item {دلیل رها کردن برخی فیلدهای هدر از محاسبه \lr{MAC}:}
    دلیل اصلی وجود {فیلدهای پویا یا متغیر در طول مسیر (\lr{Mutable Fields})} است. در پروتکل‌های شبکه، برخی از فیلدهای هدر (مانند فیلد \lr{TTL} یا \lr{Time to Live} در \lr{IP} و چک‌سام‌های لایه ۳) در هر گام (\lr{Hop}) توسط روترهای میانی تغییر می‌کنند. اگر این فیلدها وارد محاسبه \lr{MAC} شوند، به محض تغییر در اولین روتر میانی، تگ \lr{MAC} در مقصد نامعتبر شده و بسته دور انداخته می‌شود. بنابراین، طراحان مجبورند فیلدهای متغیر را از محاسبات یکپارچگی سرتاسری (\lr{End-to-End}) خارج کنند.
\end{enumerate}